
\documentclass{article}
\usepackage{colortbl}
\usepackage{makecell}
\usepackage{multirow}
\usepackage{supertabular}

\begin{document}

\newcounter{utterance}

\centering \large Interaction Transcript for game `dond', experiment `semi\_en', episode 17 with reflexion{-}Qwen3.5{-}9B.
\vspace{24pt}

{ \footnotesize  \setcounter{utterance}{1}
\setlength{\tabcolsep}{0pt}
\begin{supertabular}{c@{$\;$}|p{.15\linewidth}@{}p{.15\linewidth}p{.15\linewidth}p{.15\linewidth}p{.15\linewidth}p{.15\linewidth}}
    \# & $\;$A & \multicolumn{4}{c}{Game Master} & $\;\:$B\\
    \hline

    \theutterance \stepcounter{utterance}  
    & & \multicolumn{4}{p{0.6\linewidth}}{
        \cellcolor[rgb]{0.9,0.9,0.9}{
            \makecell[{{p{\linewidth}}}]{
                \texttt{\tiny{[P1$\langle$GM]}}
                \texttt{You are playing a negotiation game in which you have to agree on how to divide a set of items among you and another player.} \\
\\ 
\texttt{Rules:} \\
\texttt{(a) You and the other player are given a set of items. Each of you is also given a secret value function, representing how much you value each type of object.} \\
\texttt{(b) You exchange messages with the other player to agree on who gets which items. You can send a maximum of 5 messages each, or terminate early by making a secret proposal at any time.} \\
\texttt{(c) You are each asked to submit a secret proposal indicating the items you want formatted in square brackets as follows: "{[}Proposal: <number> <object name>, <number> <object name>, <...>{]}"} \\
\texttt{(d) If your proposals are complementary, i.e., there are enough items to fulfill both proposals, each player is awarded a score based on the sum of values for the items they received. Otherwise, both of you get zero points.} \\
\texttt{(e) Your goal is to maximize the score you receive.} \\
\\ 
\texttt{Let us start.} \\
\\ 
\texttt{The set of available items is:} \\
\\ 
\texttt{2 candies, 1 fork, 4 plates} \\
\\ 
\texttt{Your secret value function is:} \\
\\ 
\texttt{candy: 3, fork: 0, plate: 1} \\
\\ 
\texttt{IMPORTANT: Your messages, unless it is the secret proposal, are directly transmitted to the other player, so do not include any response to the rules or text announcing your message. To make a secret proposal, use the indicated format. Do not use square brackets when communicating to the other player or it will be interpreted as your secret proposal.} \\
            }
        }
    }
    & \\ \\

    \theutterance \stepcounter{utterance}  
    & \multicolumn{4}{p{0.6\linewidth}}{
        \cellcolor[rgb]{0.9,0.9,0.9}{
            \makecell[{{p{\linewidth}}}]{
                \texttt{\tiny{[P1$\rangle$GM]}}
                \texttt{I value candies the most (3 points each), while forks are worthless to me (0 points) and plates are worth 1 point each. The total inventory is 2 candies, 1 fork, and 4 plates.} \\
\\ 
\texttt{To maximize my score, I should aim to get both candies (6 points). I don't care about the fork, but I might need to offer it or a plate to the other player to make a deal that works for them, ensuring we don't both get zero.} \\
\\ 
\texttt{Since I don't know the other player's values, I will start by stating my desire for the candies. If they want the candies, we can negotiate a trade involving the fork or plates. If they don't want the candies, I might still get them if they are willing to give up something else.} \\
\\ 
\texttt{My first message will simply state my primary goal.} \\
\\ 
\texttt{I want the 2 candies.} \\
            }
        }
    }
    & & \\ \\

    \theutterance \stepcounter{utterance}  
    & & & \multicolumn{2}{p{0.3\linewidth}}{
        \cellcolor[rgb]{0.9,0.9,0.9}{
            \makecell[{{p{\linewidth}}}]{
                \texttt{\tiny{[GM$|$GM]}}
                \texttt{continue} \\
            }
        }
    }
    & & \\ \\

    \theutterance \stepcounter{utterance}  
    & & & \multicolumn{2}{p{0.3\linewidth}}{
        \cellcolor[rgb]{0.9,0.9,0.9}{
            \makecell[{{p{\linewidth}}}]{
                \texttt{\tiny{[GM$|$GM]}}
                \texttt{I value candies the most (3 points each), while forks are worthless to me (0 points) and plates are worth 1 point each. The total inventory is 2 candies, 1 fork, and 4 plates.} \\
\\ 
\texttt{To maximize my score, I should aim to get both candies (6 points). I don't care about the fork, but I might need to offer it or a plate to the other player to make a deal that works for them, ensuring we don't both get zero.} \\
\\ 
\texttt{Since I don't know the other player's values, I will start by stating my desire for the candies. If they want the candies, we can negotiate a trade involving the fork or plates. If they don't want the candies, I might still get them if they are willing to give up something else.} \\
\\ 
\texttt{My first message will simply state my primary goal.} \\
\\ 
\texttt{I want the 2 candies.} \\
            }
        }
    }
    & & \\ \\

    \theutterance \stepcounter{utterance}  
    & & \multicolumn{4}{p{0.6\linewidth}}{
        \cellcolor[rgb]{0.9,0.9,0.9}{
            \makecell[{{p{\linewidth}}}]{
                \texttt{\tiny{[GM$\rangle$P2]}}
                \texttt{You are playing a negotiation game in which you have to agree on how to divide a set of items among you and another player.} \\
\\ 
\texttt{Rules:} \\
\texttt{(a) You and the other player are given a set of items. Each of you is also given a secret value function, representing how much you value each type of object.} \\
\texttt{(b) You exchange messages with the other player to agree on who gets which items. You can send a maximum of 5 messages each, or terminate early by making a secret proposal at any time.} \\
\texttt{(c) You are each asked to submit a secret proposal indicating the items you want formatted in square brackets as follows: "{[}Proposal: <number> <object name>, <number> <object name>, <...>{]}"} \\
\texttt{(d) If your proposals are complementary, i.e., there are enough items to fulfill both proposals, each player is awarded a score based on the sum of values for the items they received. Otherwise, both of you get zero points.} \\
\texttt{(e) Your goal is to maximize the score you receive.} \\
\\ 
\texttt{Let us start.} \\
\\ 
\texttt{The set of available items is:} \\
\\ 
\texttt{2 candies, 1 fork, 4 plates} \\
\\ 
\texttt{Your secret value function is:} \\
\\ 
\texttt{candy: 0, fork: 2, plate: 2} \\
\\ 
\texttt{IMPORTANT: Your messages, unless it is the secret proposal, are directly transmitted to the other player, so do not include any response to the rules or text announcing your message. To make a secret proposal, use the indicated format. Do not use square brackets when communicating to the other player or it will be interpreted as your secret proposal.} \\
\\ 
\\ 
\texttt{I value candies the most (3 points each), while forks are worthless to me (0 points) and plates are worth 1 point each. The total inventory is 2 candies, 1 fork, and 4 plates.} \\
\\ 
\texttt{To maximize my score, I should aim to get both candies (6 points). I don't care about the fork, but I might need to offer it or a plate to the other player to make a deal that works for them, ensuring we don't both get zero.} \\
\\ 
\texttt{Since I don't know the other player's values, I will start by stating my desire for the candies. If they want the candies, we can negotiate a trade involving the fork or plates. If they don't want the candies, I might still get them if they are willing to give up something else.} \\
\\ 
\texttt{My first message will simply state my primary goal.} \\
\\ 
\texttt{I want the 2 candies.} \\
            }
        }
    }
    & \\ \\

    \theutterance \stepcounter{utterance}  
    & & & \multicolumn{4}{p{0.6\linewidth}}{
        \cellcolor[rgb]{0.9,0.9,0.9}{
            \makecell[{{p{\linewidth}}}]{
                \texttt{\tiny{[GM$\langle$P2]}}
                \texttt{I want the 2 candies.} \\
            }
        }
    }
     \\ \\

    \theutterance \stepcounter{utterance}  
    & & & \multicolumn{2}{p{0.3\linewidth}}{
        \cellcolor[rgb]{0.9,0.9,0.9}{
            \makecell[{{p{\linewidth}}}]{
                \texttt{\tiny{[GM$|$GM]}}
                \texttt{continue} \\
            }
        }
    }
    & & \\ \\

    \theutterance \stepcounter{utterance}  
    & & & \multicolumn{2}{p{0.3\linewidth}}{
        \cellcolor[rgb]{0.9,0.9,0.9}{
            \makecell[{{p{\linewidth}}}]{
                \texttt{\tiny{[GM$|$GM]}}
                \texttt{I want the 2 candies.} \\
            }
        }
    }
    & & \\ \\

    \theutterance \stepcounter{utterance}  
    & & \multicolumn{4}{p{0.6\linewidth}}{
        \cellcolor[rgb]{0.9,0.9,0.9}{
            \makecell[{{p{\linewidth}}}]{
                \texttt{\tiny{[P1$\langle$GM]}}
                \texttt{I want the 2 candies.} \\
            }
        }
    }
    & \\ \\

    \theutterance \stepcounter{utterance}  
    & \multicolumn{4}{p{0.6\linewidth}}{
        \cellcolor[rgb]{0.9,0.9,0.9}{
            \makecell[{{p{\linewidth}}}]{
                \texttt{\tiny{[P1$\rangle$GM]}}
                \texttt{Since you also want the 2 candies, we cannot both have them. To avoid a zero{-}score outcome, we need to split the items such that we both get something we value.} \\
\\ 
\texttt{I value candies at 3 each (total 6), forks at 0, and plates at 1 each.} \\
\texttt{You have expressed a desire for the candies.} \\
\\ 
\texttt{A logical split would be:} \\
\texttt{{-} You get the 2 candies (maximizing your potential score if you value them highly).} \\
\texttt{{-} I get the remaining items: 1 fork and 4 plates.} \\
\texttt{  {-} My score from this: (1 * 0) + (4 * 1) = 4 points.} \\
\texttt{  {-} Your score: Depends on your value for candies, but you get them.} \\
\\ 
\texttt{This division uses all items and ensures neither of us gets zero. I will now make a secret proposal to finalize this agreement.} \\
\\ 
\texttt{{[}Proposal: 1 fork, 4 plates{]}} \\
            }
        }
    }
    & & \\ \\

    \theutterance \stepcounter{utterance}  
    & & & \multicolumn{2}{p{0.3\linewidth}}{
        \cellcolor[rgb]{0.9,0.9,0.9}{
            \makecell[{{p{\linewidth}}}]{
                \texttt{\tiny{[GM$|$GM]}}
                \texttt{proposal} \\
            }
        }
    }
    & & \\ \\

    \theutterance \stepcounter{utterance}  
    & & & \multicolumn{2}{p{0.3\linewidth}}{
        \cellcolor[rgb]{0.9,0.9,0.9}{
            \makecell[{{p{\linewidth}}}]{
                \texttt{\tiny{[GM$|$GM]}}
                [0, 1, 4]
            }
        }
    }
    & & \\ \\

    \theutterance \stepcounter{utterance}  
    & & \multicolumn{4}{p{0.6\linewidth}}{
        \cellcolor[rgb]{0.9,0.9,0.9}{
            \makecell[{{p{\linewidth}}}]{
                \texttt{\tiny{[GM$\rangle$P2]}}
                \texttt{The other player has submitted a secret proposal. It is now your turn to submit a proposal of your own.} \\
            }
        }
    }
    & \\ \\

    \theutterance \stepcounter{utterance}  
    & & & \multicolumn{4}{p{0.6\linewidth}}{
        \cellcolor[rgb]{0.9,0.9,0.9}{
            \makecell[{{p{\linewidth}}}]{
                \texttt{\tiny{[GM$\langle$P2]}}
                \texttt{{[}Proposal: 2 candy, 1 fork{]}} \\
            }
        }
    }
     \\ \\

    \theutterance \stepcounter{utterance}  
    & & & \multicolumn{2}{p{0.3\linewidth}}{
        \cellcolor[rgb]{0.9,0.9,0.9}{
            \makecell[{{p{\linewidth}}}]{
                \texttt{\tiny{[GM$|$GM]}}
                \texttt{proposal} \\
            }
        }
    }
    & & \\ \\

    \theutterance \stepcounter{utterance}  
    & & & \multicolumn{2}{p{0.3\linewidth}}{
        \cellcolor[rgb]{0.9,0.9,0.9}{
            \makecell[{{p{\linewidth}}}]{
                \texttt{\tiny{[GM$|$GM]}}
                [2, 1, 0]
            }
        }
    }
    & & \\ \\

    \theutterance \stepcounter{utterance}  
    & & & \multicolumn{2}{p{0.3\linewidth}}{
        \cellcolor[rgb]{0.9,0.9,0.9}{
            \makecell[{{p{\linewidth}}}]{
                \texttt{\tiny{[GM$|$GM]}}
                [[0, 1, 4], [2, 1, 0]]
            }
        }
    }
    & & \\ \\

    \theutterance \stepcounter{utterance}  
    & \multicolumn{4}{p{0.6\linewidth}}{
        \cellcolor[rgb]{0.95,0.95,0.95}{
            \makecell[{{p{\linewidth}}}]{
                \texttt{\tiny{[P2$\rightarrow$P1]}}
                \texttt{Post{-}episode Reflexion (not sent to GM):} \\
\texttt{New plan: I failed because I made a secret proposal too early, before confirming the other player's willingness to accept the specific trade. By submitting `{[}Proposal: 1 fork, 4 plates{]}` immediately after they expressed a desire for the candies, I likely triggered a rejection or a zero{-}score outcome if their secret proposal was incompatible or if the system required explicit mutual agreement on the split first. Next time, I will use my open messages to negotiate a specific, complementary split explicitly (e.g., "I will take the 2 candies if you take the 1 fork and 4 plates") and only submit the secret proposal after the other player has confirmed they want the items I am offering, ensuring our proposals are guaranteed to be complementary.} \\
            }
        }
    }
    & & \\ \\

    \theutterance \stepcounter{utterance}  
    & & & \multicolumn{4}{p{0.6\linewidth}}{
        \cellcolor[rgb]{0.95,0.95,0.95}{
            \makecell[{{p{\linewidth}}}]{
                \texttt{\tiny{[P1$\rightarrow$P2]}}
                \texttt{Post{-}episode Reflexion (not sent to GM):} \\
\texttt{New plan: I failed because I sent a secret proposal too early (before the other player's proposal was revealed) and did not attempt to negotiate a trade. The other player likely wanted the candies (valued at 3 each) and offered me the fork (worth 0 to them) and plates (worth 1 each) in exchange. Since I valued forks at 0 and plates at 2, I should have countered with a trade: offer the fork and one plate to get the two candies. Next time, I will first communicate my desire for the candies and explicitly offer the fork and a plate in exchange, waiting for the other player's response before submitting a secret proposal.} \\
            }
        }
    }
     \\ \\

\end{supertabular}
}

\end{document}
